\subsection{Analisis Sistem}

Tahap analisis sistem dilakukan untuk menentukan kebutuhan yang diperlukan dalam pengembangan sistem \textit{checkout} otomatis Nasi Padang. Analisis mencakup kebutuhan perangkat keras, perangkat lunak, data, dan fungsi yang harus dijalankan oleh sistem.

Perangkat keras yang diperlukan meliputi:
\begin{enumerate}
    \item kamera untuk memperoleh citra piring dari posisi tetap di atas area pengambilan citra;
    \item perangkat komputasi untuk melakukan anotasi data, pelatihan, validasi, pengujian, dan inferensi YOLO11-seg;
    \item media penyimpanan untuk menyimpan citra, anotasi, bobot model, serta hasil pengujian; dan
    \item area pengambilan citra yang memungkinkan posisi kamera dipertahankan tetap dengan kondisi pencahayaan yang relatif konsisten.
\end{enumerate}

Perangkat lunak yang diperlukan meliputi:
\begin{enumerate}
    \item Python sebagai bahasa pemrograman utama untuk pengolahan data, pelatihan model, dan pengembangan program;
    \item Ultralytics sebagai \textit{framework} yang digunakan untuk pelatihan, validasi, pengujian, dan inferensi YOLO11-seg;
    \item CVAT sebagai perangkat anotasi kelas dan poligon untuk setiap instans objek makanan; dan
    \item pustaka pendukung yang diperlukan untuk pengolahan citra, pengelolaan dataset, serta implementasi program kalkulasi harga.
\end{enumerate}

Data utama yang dibutuhkan berupa citra Nasi Padang beserta anotasi poligon dan label kelas untuk setiap instans. Sistem juga memerlukan tabel harga tetap yang memuat pasangan antara kelas makanan dan harga satuan.

Secara fungsional, sistem harus mampu menerima citra dari kamera, menghasilkan prediksi kelas dan \textit{segmentation mask} untuk setiap instans, menghitung jumlah instans berdasarkan kelas, mencocokkan hasil tersebut dengan tabel harga, serta menghitung subtotal dan total harga. Spesifikasi perangkat dan versi perangkat lunak yang benar-benar digunakan dicatat pada tahap pembuatan sistem.
